\documentclass[11pt]{article}
\usepackage[utf8]{inputenc}
\usepackage[margin=1in]{geometry}
\usepackage{amsmath}
\usepackage{graphicx}
\usepackage{booktabs}
\usepackage{hyperref}
\usepackage[round]{natbib}
\usepackage{float}

\title{Modulation of Alpha Oscillations by Attention Is Predicted by Hemispheric Asymmetry of Subcortical Regions}
\author{Author A$^{1}$, Author B$^{1}$, Author C$^{1,2}$\\
\small $^{1}$Centre for Human Brain Health, University of Birmingham, UK\\
\small $^{2}$Department of Psychology, University of Birmingham, UK}
\date{}

\begin{document}
\maketitle

\begin{abstract}
Spatial attention modulates posterior alpha-band (8--12~Hz) oscillations, with power decreasing contralateral and increasing ipsilateral to the attended hemifield. Although cortical contributions from dorsal and ventral attention networks are well characterised, the role of subcortical structures in shaping individual differences in alpha lateralisation remains poorly understood. Here, we combined magnetoencephalography (MEG) with structural MRI volumetrics in 33 right-handed participants who performed a cued spatial attention change-detection task under a $2 \times 2$ factorial design crossing target perceptual load (high vs.\ low) with distractor salience (high vs.\ low). We computed a hemispheric lateralisation modulation index for alpha power (HLM$_\alpha$) and lateralisation volumes (LV) for five subcortical nuclei extracted via FreeSurfer. A generalized linear model (GLM) selected by AIC/BIC identified thalamus, caudate nucleus, and globus pallidus LV as the best predictors of HLM$_\alpha$ ($p = 0.0007$). Condition-specific multivariate regression ($p = 0.001$) revealed a striking dissociation: thalamus LV predicted alpha lateralisation under low task demands, globus pallidus LV under intermediate demands, and caudate nucleus LV under the highest demands. These results establish, for the first time in humans, a link between structural hemispheric asymmetry of subcortical nuclei and functional alpha-band lateralisation during spatial attention, implicating a load-dependent hierarchy of subcortical contributions to attentional oscillatory control.
\end{abstract}

\section{Introduction}

Spatial attention profoundly shapes sensory processing by biasing neural activity toward behaviourally relevant locations in the visual field. A hallmark electrophysiological signature of this process is the lateralisation of posterior alpha-band oscillations (8--12~Hz): alpha power decreases over cortical regions contralateral to the attended hemifield and increases ipsilaterally \citep{worden2000anticipatory, thut2006alpha, foxe2011role}. This pattern is widely interpreted as reflecting the gating of information flow, whereby alpha suppression facilitates processing of attended stimuli while alpha enhancement inhibits processing of distractors \citep{klimesch2007eeg, jensen2002oscillations}.

Extensive research has delineated the cortical networks mediating spatial attention, including the dorsal attention network (frontal eye fields, intraparietal sulcus) and the ventral attention network \citep{corbetta2002control, petersen2012attention, kastner2000mechanisms}. However, subcortical contributions to attentional alpha modulation remain poorly understood. Animal studies have established that the pulvinar nucleus of the thalamus regulates information transmission between cortical areas based on attentional demands \citep{saalmann2012neural, mcalonan2008guarding}, and that basal ganglia nuclei participate in attentional selection and gating \citep{desimone1995neural, lanciego2012functional}. In humans, lesion evidence implicates the pulvinar in attentional selection \citep{snow2009impaired}, and the thalamus more broadly in cognitive control \citep{halassa2014thalamic, saalmann2011cognitive}. Yet direct evidence linking subcortical structure to oscillatory attentional mechanisms in the healthy human brain is scarce.

A major methodological barrier is that MEG and EEG cannot reliably detect deep-source activity from the thalamus and basal ganglia due to their distance from sensors and low signal-to-noise ratio. This limitation has left a gap between animal electrophysiology, which can record directly from these structures, and human neuroimaging, which typically relies on haemodynamic measures lacking temporal resolution. Here, we adopt a novel indirect approach: rather than attempting to measure subcortical oscillatory activity, we test whether structural hemispheric asymmetries of subcortical nuclei predict individual differences in functional alpha lateralisation during spatial attention.

Furthermore, the relationship between perceptual load and subcortical engagement during attention is largely unexplored. Load theory \citep{lavie2005distracted} predicts that distractor interference depends on available processing capacity, but how subcortical structures contribute to this capacity-dependent modulation of alpha oscillations has not been examined. Given that the thalamus, caudate nucleus, and globus pallidus have distinct connectivity profiles and functional roles, we hypothesised that they might differentially predict alpha lateralisation under varying levels of task demand.

To test these hypotheses, we re-analysed a previously collected dataset in which 33 participants performed a cued spatial attention change-detection task while MEG and structural MRI were recorded. The task employed a $2 \times 2$ factorial design crossing target perceptual load with distractor salience, generating four conditions of varying attentional demand. We extracted bilateral volumes of five subcortical structures using FreeSurfer segmentation, computed lateralisation volumes, and used GLM model selection to identify which subcortical asymmetries best predicted individual variation in alpha lateralisation across conditions.

\section{Results}

\subsection{Behavioural Performance and Alpha Modulation}

Participants detected gaze-shift direction in the cued face across all four conditions. Performance varied as expected, with higher perceptual load reducing accuracy and higher distractor salience increasing interference. The $2 \times 2$ factorial design (Table~\ref{tab:conditions}) produced systematic variation in task demands.

\begin{table}[H]
\centering
\caption{Experimental conditions of the $2 \times 2$ factorial design.}
\label{tab:conditions}
\begin{tabular}{@{}lllll@{}}
\toprule
Condition & Target Load & Distractor Salience & Target & Distractor \\
\midrule
1 & Low & Low & Clear face & Noisy face \\
2 & High & Low & Noisy face & Noisy face \\
3 & Low & High & Clear face & Clear face \\
4 & High & High & Noisy face & Clear face \\
\bottomrule
\end{tabular}
\end{table}

Time--frequency analysis of MEG data during the pre-target interval ($-850$ to $0$~ms) confirmed robust alpha-band modulation consistent with spatial attention. Contralateral sensors showed decreased alpha power while ipsilateral sensors showed relative increases. The modulation index MI$_\alpha$ increased gradually during the delay interval until target onset, reflecting progressive attentional engagement. Five symmetric sensor clusters per hemisphere showing the strongest modulation were selected as regions of interest (ROIs). The distribution of HLM$_\alpha$ across participants showed no systematic hemispheric bias at the group level ($t(32) = 0.87$, $p = 0.39$, two-tailed $t$-test), indicating substantial individual variability.

\subsection{Subcortical Volume Lateralisation}

Bilateral volumes of five subcortical structures were extracted from structural MRI using FreeSurfer \citep{fischl2002whole}. Lateralisation volume (LV) was computed as LV $= (\text{Left} - \text{Right}) / (\text{Left} + \text{Right})$ for each structure (Table~\ref{tab:lateralisation}).

\begin{table}[H]
\centering
\caption{Subcortical volume lateralisation across 33 participants.}
\label{tab:lateralisation}
\begin{tabular}{@{}lccc@{}}
\toprule
Structure & Mean LV ($\times 10^{-2}$) & $t$-statistic & $p$-value \\
\midrule
Caudate nucleus & $-1.42$ & $-2.43$ & \textbf{0.021} \\
Putamen & $+0.31$ & $0.58$ & 0.566 \\
Nucleus accumbens & $-0.67$ & $-1.14$ & 0.263 \\
Globus pallidus & $+0.48$ & $0.91$ & 0.370 \\
Thalamus & $+0.79$ & $1.73$ & 0.093 \\
\bottomrule
\end{tabular}
\end{table}

Only the caudate nucleus showed significant lateralisation, being right-lateralised ($p = 0.021$, one-sample $t$-test against zero). The thalamus showed a non-significant trend toward left lateralisation ($p = 0.093$). Putamen, nucleus accumbens, and globus pallidus showed no significant lateralisation.

\subsection{GLM Model Selection}

To determine which subcortical structures best predict HLM$_\alpha$, we evaluated all possible combinations of the five subcortical LV regressors using both AIC and BIC \citep{burnham2002practical}. Table~\ref{tab:model_selection} summarises the model comparison.

\begin{table}[H]
\centering
\caption{GLM model selection: AIC and BIC for all regressor combinations. The best model is indicated in bold.}
\label{tab:model_selection}
\begin{tabular}{@{}lccc@{}}
\toprule
Model (Regressors) & AIC & $\Delta$AIC & $\Delta$BIC \\
\midrule
Thalamus only & $-48.3$ & $12.1$ & $11.8$ \\
Caudate only & $-51.7$ & $8.7$ & $8.4$ \\
Globus pallidus only & $-49.1$ & $11.3$ & $11.0$ \\
Putamen only & $-44.2$ & $16.2$ & $15.9$ \\
Nucleus accumbens only & $-43.8$ & $16.6$ & $16.3$ \\
Thalamus + Caudate & $-55.2$ & $5.2$ & $5.8$ \\
Thalamus + GP & $-53.9$ & $6.5$ & $7.1$ \\
Caudate + GP & $-54.1$ & $6.3$ & $6.9$ \\
\textbf{Thalamus + Caudate + GP} & $\mathbf{-60.4}$ & $\mathbf{0.0}$ & $\mathbf{0.0}$ \\
All 5 structures & $-57.1$ & $3.3$ & $6.2$ \\
\bottomrule
\end{tabular}
\end{table}

The three-regressor model including thalamus, caudate nucleus, and globus pallidus LV yielded the lowest AIC ($-60.4$) and BIC values, and was therefore selected as the best model. The full five-structure model did not improve fit, indicating that putamen and nucleus accumbens do not contribute additional predictive power.

\subsection{Best GLM Model: Pooled Analysis}

The selected three-regressor GLM significantly predicted HLM$_\alpha$ compared to the null model ($F(3, 29) = 7.82$, $p = 0.0007$). All three subcortical structures showed significant positive beta coefficients (Table~\ref{tab:betas}).

\begin{table}[H]
\centering
\caption{Beta coefficients for the best GLM model predicting HLM$_\alpha$ (pooled across conditions).}
\label{tab:betas}
\begin{tabular}{@{}lcccc@{}}
\toprule
Regressor & $\beta$ & SE & $t$-statistic & $p$-value \\
\midrule
Thalamus LV & $0.347$ & $0.128$ & $2.71$ & \textbf{0.011} \\
Caudate LV & $0.312$ & $0.119$ & $2.62$ & \textbf{0.014} \\
Globus pallidus LV & $0.298$ & $0.131$ & $2.27$ & \textbf{0.031} \\
\bottomrule
\end{tabular}
\end{table}

Each structure showed a positive relationship with HLM$_\alpha$: larger left-relative-to-right subcortical volume was associated with stronger left-hemisphere alpha modulation relative to right. The model explained a substantial proportion of variance in individual alpha lateralisation differences.

\subsection{Condition-Specific Multivariate Regression}

To test whether subcortical contributions vary with task demands, we performed a multivariate regression predicting HLM$_\alpha$ separately in each of the four conditions using the same three regressors. The overall multivariate model was significant ($F(12, 84) = 2.94$, $p = 0.001$). Condition-specific beta coefficients revealed a striking dissociation (Table~\ref{tab:condition_betas}).

\begin{table}[H]
\centering
\caption{Condition-specific beta coefficients from multivariate regression. Significant predictors ($p < 0.05$) in bold.}
\label{tab:condition_betas}
\begin{tabular}{@{}lcccccc@{}}
\toprule
Condition & $\beta_{\text{Thal}}$ & $p_{\text{Thal}}$ & $\beta_{\text{Caud}}$ & $p_{\text{Caud}}$ & $\beta_{\text{GP}}$ & $p_{\text{GP}}$ \\
\midrule
1: Low/Low & \textbf{0.491} & \textbf{0.004} & $0.127$ & $0.312$ & $0.083$ & $0.521$ \\
2: High/Low & $0.168$ & $0.187$ & $0.094$ & $0.467$ & \textbf{0.384} & \textbf{0.018} \\
3: Low/High & $0.142$ & $0.244$ & $0.112$ & $0.389$ & \textbf{0.401} & \textbf{0.013} \\
4: High/High & $0.079$ & $0.541$ & \textbf{0.467} & \textbf{0.003} & $0.131$ & $0.298$ \\
\bottomrule
\end{tabular}
\end{table}

Under the lowest task demands (Condition 1: clear target, noisy distractor), thalamus LV was the sole significant predictor ($\beta = 0.491$, $p = 0.004$). Under intermediate demands (Conditions 2 and 3: either high load or high salience), globus pallidus LV became the dominant predictor ($\beta = 0.384$, $p = 0.018$; $\beta = 0.401$, $p = 0.013$). Under the highest demands (Condition 4: noisy target, clear distractor), caudate nucleus LV was the dominant predictor ($\beta = 0.467$, $p = 0.003$). This dissociation suggests a load-dependent hierarchy of subcortical contributions to alpha lateralisation.

\subsection{Summary of Subcortical Contributions by Task Demand}

Table~\ref{tab:demand_summary} summarises the dissociation across demand levels.

\begin{table}[H]
\centering
\caption{Summary of dominant subcortical predictor by task demand level.}
\label{tab:demand_summary}
\begin{tabular}{@{}lll@{}}
\toprule
Demand Level & Dominant Predictor & Condition Description \\
\midrule
Lowest & Thalamus & Clear target, noisy distractor \\
Intermediate & Globus pallidus & High load or high salience \\
Highest & Caudate nucleus & Noisy target, clear distractor \\
\bottomrule
\end{tabular}
\end{table}

\subsection{Supplementary Analyses: RIFT and Behavioural Asymmetry}

We additionally tested whether the three-regressor model predicted hemispheric lateralisation of rapid invisible frequency tagging (RIFT) signals and behavioural asymmetry. Bayes factors (BF) were computed for correlations between each subcortical structure's LV and HLM$_\text{RIFT}$ as well as behavioural asymmetry (Table~\ref{tab:bayes}).

\begin{table}[H]
\centering
\caption{Bayes factors for correlations between subcortical LV and supplementary measures. BF $> 3$ indicates moderate evidence for the alternative hypothesis; BF $< 1/3$ indicates moderate evidence for the null.}
\label{tab:bayes}
\begin{tabular}{@{}lcc@{}}
\toprule
Structure & BF (HLM$_\text{RIFT}$) & BF (Behavioural Asymmetry) \\
\midrule
Thalamus & $2.14$ & $0.87$ \\
Caudate nucleus & $1.67$ & $0.54$ \\
Globus pallidus & $1.43$ & $0.61$ \\
Putamen & $0.38$ & $0.29$ \\
Nucleus accumbens & $0.31$ & $0.24$ \\
\bottomrule
\end{tabular}
\end{table}

Bayes factors for the three model regressors (thalamus, caudate, globus pallidus) were in the anecdotal-to-moderate range for HLM$_\text{RIFT}$, suggesting a trend consistent with the alpha findings. For behavioural asymmetry, evidence was largely inconclusive, suggesting that structural subcortical asymmetry primarily influences neural rather than behavioural lateralisation.

\section{Discussion}

This study demonstrates, for the first time in humans, that structural hemispheric asymmetry of subcortical nuclei predicts individual differences in attention-related alpha oscillation lateralisation. Using a novel indirect approach that combines structural MRI volumetrics with MEG-measured alpha power, we bypassed the fundamental limitation of directly measuring subcortical oscillatory activity with non-invasive methods. Three key findings emerged.

First, among five subcortical structures examined, the thalamus, caudate nucleus, and globus pallidus collectively provided the best prediction of alpha lateralisation (GLM $p = 0.0007$), while putamen and nucleus accumbens did not contribute additional predictive power. This selectivity is consistent with the known involvement of the thalamus (particularly the pulvinar) in attentional gating \citep{saalmann2012neural, mcalonan2008guarding} and of basal ganglia in selection mechanisms \citep{lanciego2012functional, desimone1995neural}. The positive direction of all three beta coefficients indicates that individuals with relatively larger left subcortical volumes exhibit stronger left-hemisphere alpha modulation, suggesting a structural basis for functional lateralisation.

Second, condition-specific analyses revealed a striking dissociation: the thalamus dominated under low demands, the globus pallidus under intermediate demands, and the caudate nucleus under the highest demands. This pattern aligns with a hierarchical model in which different subcortical structures are recruited as attentional demands escalate. The thalamus, with its established role in sensory gating \citep{saalmann2011cognitive}, may provide a default attentional bias that suffices for easy discriminations. The globus pallidus, a key output nucleus of the basal ganglia, may be engaged when either perceptual load or distractor salience increases but not both simultaneously. The caudate nucleus, which receives extensive prefrontal projections and participates in cognitive control \citep{lanciego2012functional}, may be recruited only when the most demanding combination of high load and high salience distractor requires top-down override.

Third, the observation that caudate nucleus volume is significantly right-lateralised ($p = 0.021$) while the thalamus trends toward left lateralisation introduces an interesting structural asymmetry that may constrain individual variation in attentional alpha lateralisation. Such volumetric asymmetries have been documented in healthy populations but their functional significance has been unclear. Our results suggest they may reflect individual differences in the balance of subcortical contributions to oscillatory attentional control.

Several limitations should be noted. First, the correlational nature of our volumetric--functional approach cannot establish causality. Second, FreeSurfer volumetrics provide aggregate measures that do not distinguish among thalamic or basal ganglia sub-nuclei. Third, the sample size of 33 participants, while adequate for the regression models employed, limits power for detecting smaller effects. Fourth, the study used a re-analysis design, so the task was not specifically optimised for the current hypotheses.

The clinical implications of these findings merit consideration. Structural atrophy in subcortical regions is a hallmark of several neurological disorders including Parkinson's disease, Huntington's disease, and Alzheimer's disease. If subcortical asymmetry shapes oscillatory attention, then disease-related volume changes could produce measurable alterations in alpha lateralisation, potentially serving as early biomarkers. Future longitudinal studies in clinical populations could test this prediction.

In conclusion, subcortical volumetric asymmetry predicts the degree and lateralisation of attention-related alpha modulation, with a load-dependent hierarchy whereby the thalamus, globus pallidus, and caudate nucleus are differentially recruited as task demands increase. These findings open a new line of investigation into the subcortical structural bases of oscillatory attentional control.

\section{Methods}

\subsection{Participants}

Thirty-five right-handed healthy volunteers were recruited; two were excluded due to consent withdrawal or MRI quality issues, yielding a final sample of 33 (25 female, mean age $24 \pm 5.7$ years). All had normal or corrected-to-normal vision and were compensated at 15~GBP per hour. The study was approved by the STEM Ethics Committee at the University of Birmingham.

\subsection{Experimental Design}

Participants performed a cued spatial attention change-detection task across 512 trials (2 blocks $\times$ 256 trials, $\sim$45~min). On each trial, participants fixated centrally while face stimuli (8 identities, 8$^\circ$ diameter) appeared at 7$^\circ$ eccentricity in both hemifields. A directional cue (350~ms) indicated the target side. After a variable delay (1000--2000~ms), the gaze of each face shifted (150~ms), and participants reported the target gaze direction within 1000~ms. A $2 \times 2$ factorial manipulation crossed target load (clear vs.\ noisy face; 50\% pixel noise) with distractor salience (clear vs.\ noisy distractor), yielding four conditions (Table~\ref{tab:conditions}). Stimulus presentation used Psychophysics Toolbox 3.0.11 \citep{brainard1997psychophysics} on MATLAB/Windows~10.

\subsection{MEG Acquisition and Analysis}

MEG was recorded using a 306-channel Elekta/MEGIN Neuromag system (204 planar gradiometers, 102 magnetometers) at 1000~Hz with online bandpass filtering (0.1--330~Hz). Five head position indicator coils and Polhemus digitisation tracked head position. Alpha power modulation was quantified in the $-850$ to $0$~ms pre-target interval. A modulation index MI$_\alpha$ reflecting the relative difference between attend-left and attend-right trials was computed at each sensor. Five symmetric sensor clusters per hemisphere showing strongest modulation were selected as ROIs (10 sensors total). The hemispheric lateralisation modulation HLM$_\alpha$ was computed as the difference between left and right hemisphere modulation indices, normalised by their sum, to capture individual differences in hemispheric alpha modulation balance.

\subsection{Structural MRI and Volumetric Analysis}

Structural MRI was acquired at 3T. Automated segmentation using FreeSurfer \citep{fischl2002whole} extracted bilateral volumes of the thalamus, caudate nucleus, putamen, globus pallidus, and nucleus accumbens. Lateralisation volume was computed as LV $= (\text{Left} - \text{Right}) / (\text{Left} + \text{Right})$ for each structure. One-sample $t$-tests assessed whether LV differed from zero.

\subsection{Statistical Modelling}

A GLM tested whether subcortical LV values predicted HLM$_\alpha$ (pooled across conditions). Model selection used AIC and BIC across all $2^5 - 1 = 31$ possible regressor combinations \citep{burnham2002practical}. A multivariate regression then examined how the selected regressors predicted HLM$_\alpha$ separately in each condition. Bayes factors were computed for correlations between subcortical structures and both HLM$_\text{RIFT}$ and behavioural asymmetry, comparing alternative versus null hypotheses \citep{zhigalov2019alpha, lozanosoldevilla2024rift}.

\section{Conclusion}

We established that hemispheric volumetric asymmetry of three subcortical nuclei---thalamus, caudate nucleus, and globus pallidus---significantly predicts individual variation in attention-driven alpha oscillation lateralisation. A load-dependent dissociation reveals that distinct subcortical structures dominate under different attentional demands, suggesting a hierarchical subcortical architecture for oscillatory control of spatial attention. These findings bridge a critical gap between animal electrophysiology demonstrating subcortical involvement in attention and the scarcity of corresponding human evidence, offering a structural framework for understanding individual differences in oscillatory attentional mechanisms.

\bibliographystyle{plainnat}
\bibliography{references}

\end{document}
